\author{Brett Reynolds} %use this field for editors as well
\title{言語の風景：Language Landscapes}
\subtitle{英語教師のための英語のしくみガイド【日本語版】}
% \ISBNdigital{}
% \ISBNhardcover{}
% \BookDOI{}
\typesetter{Brett Reynolds}
% \proofreader{}
% \lsCoverTitleSizes{51.5pt}{17pt}
\renewcommand{\lsSeries}{tbls}
%\newcommand{\lsSeriesTitle}{Textbooks in Language Sciences}
%\newcommand{\lsSeriesColor}{lsYellow}
\renewcommand{\lsSeriesNumber}{16}
\BackBody{言語は暗記するためのリストではなく、探求すべき風景です。

優れたフィールド・ナチュラリストは、単に生き物に名前を付けるだけでなく、それらがどのように振る舞い、どこに群生し、なぜ変化するのかを観察します。本書では、英語を同じように読み解くことを求めています。つまり、ルールブックがどうあるべきかと言うことではなく、話者が実際に何をしているかに注目するのです。理論的枠組みは\textit{The Cambridge Grammar of the English Language}（ケンブリッジ英文法）に基づいており、理解しやすいようにシンプルにされていますが、決してレベルを下げているわけではありません。

各章は、音声や単語から句、節、そして談話へと進み、英語の明確な例が他の言語のワンシーンとともに提示されています。それぞれの章で同じプロセスを実践します：パターンを観察し、証拠と照らし合わせて検証し、どこでそれが破綻するかに気づき、その破綻が何を物語っているのかを問いかけます。最後まで読み終える頃には、自らの主張を裏付けることができる実用的な文法知識が身についていることでしょう。それは、教科書を評価したり、学習者の文がなぜ少し不自然なのかを説明したり、疑問を投げかけられたときに自分の説明を弁護したりするために使える知識です。

あなたが教員養成の過程にあっても、すでに教室で教えていても、得られるものは同じです。根拠のないルールに頼るのをやめ、あなた自身で言語システムを理解できるようになるのです。}